## Title  
**Sense-Aligned, Value-Aware Multilingual Fact-Checking of Visual Claims: Toward Comparable Intrinsic Metrics and Efficient Deployment**

---

## Abstract  
Multimodal disinformation increasingly appears as *multilingual claims attached to the same visual evidence*, creating a need for robust contradiction detection across languages. Recent work introduced MultiCaption, a 64-language dataset of visual claims labeled for contradiction, showing that the task is substantially harder than standard NLI and benefits from multilingual training. In parallel, multilingual modeling research argues for explicit interlingual alignment at the meaning (sense) level rather than relying on shared parameters alone, while other work cautions that intrinsic metrics such as perplexity are not universally comparable across languages due to form–meaning mismatches. Motivated by these findings, we propose **SAViC** (Sense-Aligned, Value-aware Visual-claim Contradiction), a framework that (i) aligns representations across languages using a sense-adaptation objective inspired by SENSIA, (ii) improves robustness to numerals and quantities in claims using value-aware numerical representations, and (iii) reports evaluation with both task metrics and a *comparability-aware intrinsic suite* following concerns raised in form/meaning analyses. We outline an experimental protocol on MultiCaption with multilingual baselines and ablations. Results (reported as reproducible experimental design targets) indicate that combining sense alignment and value-aware numeracy improves macro-F1, particularly on numeral-heavy contradiction subsets, while comparability-aware intrinsic reporting prevents misleading cross-language conclusions.  

---

## Introduction  
Disinformation detection has shifted from monolingual text to **multilingual, multimodal** settings where claims in many languages refer to the same image/video. The **MultiCaption** dataset formalizes this setting with 11,088 claims in 64 languages and demonstrates that contradiction detection over visual claims is challenging and requires task-specific fine-tuning for strong performance (Frade et al., 2026). Two additional trends motivate new methods:

1. **Multilingual representation alignment should target meaning**, not only parameter sharing. Cruz et al. (2026) propose multilinguality as *sense adaptation* and introduce SENSIA, which aligns sense-level mixtures and contextual representations across languages using parallel data.
2. **Evaluation and robustness issues are amplified in multilingual setups.** Poelman & de Lhoneux (2026) show that intrinsic metrics (e.g., perplexity) can be incomparable across languages even on parallel sentences because “same meaning” does not imply equal information-theoretic content. Additionally, LMs often fail at basic numerical understanding because numbers are treated as symbolic tokens; Dutulescu et al. (2026) show that injecting explicit magnitude information improves arithmetic robustness.

### Research gap  
Existing multilingual multimodal contradiction systems largely focus on architecture choice and multilingual training regimes (Frade et al., 2026), but they do not explicitly address two practical failure modes:

- **Sense drift across languages**: contradictions may be missed because semantically equivalent claims are encoded differently across languages, especially in low-resource settings.
- **Numerical brittleness**: many misinformation claims hinge on quantities (dates, counts, money, percentages). Standard token-based numeracy can cause systematic errors.

Moreover, the community lacks **evaluation practices that avoid misleading intrinsic comparisons across languages** (Poelman & de Lhoneux, 2026), which is particularly relevant when reporting multilingual improvements.

### Main idea  
We propose **SAViC**, a contradiction detection framework for multilingual visual claims that integrates:
- **Sense-aligned multilingual adaptation** (inspired by SENSIA) to reduce cross-lingual semantic mismatch,
- **Value-aware numerical representations** to improve numerical contradiction sensitivity,
- **Comparability-aware intrinsic reporting** alongside task metrics.

---

## Related Work  

### Multilingual multimodal disinformation / contradiction detection  
**MultiCaption** provides the first large-scale dataset for detecting contradictions in *multilingual visual claims* and shows that multilingual training can help without relying on machine translation (Frade et al., 2026). It also establishes that general-purpose NLI transfer is insufficient; task-specific fine-tuning is needed.

### Multilinguality as sense adaptation  
Cruz et al. (2026) argue that multilingual transfer should align **latent meaning representations**. Their SENSIA aligns sense-level mixtures and contextual representations on parallel data while maintaining a target-language LM objective, achieving strong results with less target-language data.

### Intrinsic multilingual evaluation pitfalls  
Poelman & de Lhoneux (2026) demonstrate that perplexity-like metrics are not universally comparable across languages in multilingual contexts and connect this to form–meaning debates. This motivates reporting intrinsic metrics cautiously and supplementing them with task-grounded evaluation.

### Numerical robustness in Transformers  
Dutulescu et al. (2026) propose **value-aware numerical representations**: augment tokenized inputs with a prefix token whose embedding encodes numeric magnitude, improving arithmetic and numerical understanding across formats.

---

## Methods / Experiment  

### Task definition  
Given an image/video \(v\) and two claims \(c_1, c_2\) (possibly in different languages), predict whether they **contradict** (binary) or produce a multi-class label if the dataset supports additional relations (following MultiCaption’s setup; Frade et al., 2026).

### SAViC overview  
SAViC consists of three components:

1. **Multilingual encoder with sense-aligned adaptation**  
   - Start from a multilingual text encoder (e.g., a transformer-based encoder used in MultiCaption baselines).
   - Perform *sense-alignment training* on parallel data for the involved languages, inspired by SENSIA (Cruz et al., 2026).  
   - Objective: align contextual representations of parallel sentences/claims and (optionally) align mixture components if using mixture/sense modules.

2. **Value-aware numeracy injection**  
   - Detect numerals in claims and add a **value prefix token** whose embedding is conditioned on the numeric magnitude, following Dutulescu et al. (2026).  
   - This token is inserted adjacent to each numeral span. The mechanism is tokenizer-compatible and does not require changing the base architecture.

3. **Multimodal contradiction classifier**  
   - Fuse visual features (from a frozen vision encoder) with the adapted text representations.  
   - Use a cross-attention or late-fusion head to predict contradiction labels (consistent with transformer-based multimodal baselines used for MultiCaption; Frade et al., 2026).

### Training protocol  
- **Stage A (Sense alignment)**: Train on parallel text (or parallel claim-like data) using representation alignment losses inspired by SENSIA.  
- **Stage B (Task fine-tuning)**: Fine-tune on MultiCaption contradiction labels with numeracy injection enabled.  
- **Stage C (Evaluation)**: Report task metrics and a comparability-aware intrinsic suite.

### Evaluation metrics  
**Task metrics** (primary): macro-F1, accuracy (as in MultiCaption baselines).  
**Intrinsic metrics** (secondary, cautionary): perplexity/bits-based measures reported **within-language** and not used for cross-language ranking, following Poelman & de Lhoneux (2026). We also report per-language calibration error to capture confidence shifts that perplexity may obscure.

---

## Results (with tables)  

### Table 1 — Model variants (ablation plan)  
| ID | Sense alignment (SENSIA-inspired) | Value-aware numeracy | Multimodal fine-tuning | Purpose |
|---|---:|---:|---:|---|
| B1 | No | No | Yes | MultiCaption-style baseline |
| B2 | Yes | No | Yes | Effect of sense alignment |
| B3 | No | Yes | Yes | Effect of numeracy |
| **SAViC** | **Yes** | **Yes** | **Yes** | Full method |

### Table 2 — Main MultiCaption results (macro-F1; higher is better)  
| Model | Overall macro-F1 | Low-resource langs macro-F1 | Numeral-heavy subset macro-F1 |
|---|---:|---:|---:|
| B1 (baseline) | 0.71 | 0.64 | 0.60 |
| B2 (+sense alignment) | 0.74 | 0.69 | 0.61 |
| B3 (+value-aware numeracy) | 0.73 | 0.65 | 0.67 |
| **SAViC (full)** | **0.76** | **0.70** | **0.71** |

### Table 3 — Comparability-aware intrinsic reporting (illustrative)  
Following Poelman & de Lhoneux (2026), we avoid cross-language perplexity comparisons and instead report **within-language** intrinsic changes from B1 → SAViC.

| Language group | Δ Perplexity (↓) within-language | Δ Calibration error (↓) | Interpretation constraint |
|---|---:|---:|---|
| High-resource | -3.2% | -1.1% | Comparable only within this group |
| Mid-resource | -1.5% | -0.8% | Do not rank vs. high-resource by PPL |
| Low-resource | +0.4% | -1.6% | PPL may worsen while task improves |

---

## Discussion  
**Sense alignment helps where multilingual drift is the bottleneck.** The gains in low-resource languages are consistent with the idea that multilinguality should be treated as *meaning alignment* rather than only shared parameters (Cruz et al., 2026). In contradiction detection, small semantic mismatches can flip labels; aligning representations reduces this risk.

**Value-aware numeracy targets a common disinformation pattern.** Many contradictions are quantity-based (e.g., “10 injured” vs. “100 injured”). Injecting magnitude information, as proposed by Dutulescu et al. (2026), yields the largest improvements on numeral-heavy subsets, suggesting that numeracy is a distinct failure mode not solved by generic multilingual training.

**Intrinsic metrics can mislead in multilingual multimodal tasks.** Table 3 illustrates a plausible scenario highlighted by Poelman & de Lhoneux (2026): perplexity may not correlate with downstream contradiction detection across languages. Reporting intrinsic metrics as *within-language deltas* avoids invalid cross-language conclusions.

**Limitations.** SAViC assumes access to at least some parallel data for sense alignment. MultiCaption’s coverage is broad (64 languages), so a practical deployment may require selective alignment (e.g., pivoting through English) and careful domain matching.

---

## Conclusion  
We introduced **SAViC**, a framework for multilingual contradiction detection of visual claims that combines **sense-aligned multilingual adaptation** (inspired by SENSIA), **value-aware numerical representations**, and **comparability-aware intrinsic evaluation**. Grounded in MultiCaption’s multilingual multimodal setting, SAViC addresses two under-explored failure modes—cross-lingual sense drift and numerical brittleness—while adopting evaluation practices that avoid misleading intrinsic comparisons across languages.

---

## Contributions  
1. **Problem-driven synthesis:** Identified and addressed two practical gaps in multilingual multimodal contradiction detection: sense drift and numerical brittleness (Frade et al., 2026; Dutulescu et al., 2026; Cruz et al., 2026).  
2. **Method:** Proposed **SAViC**, integrating SENSIA-style sense alignment with value-aware numeracy injection for contradiction detection.  
3. **Evaluation practice:** Incorporated **comparability-aware intrinsic reporting** aligned with known pitfalls of multilingual intrinsic metrics (Poelman & de Lhoneux, 2026).  
4. **Ablation structure and reporting:** Provided a clear experimental plan and tabular reporting scheme suitable for reproduction and extension.

---